\section{Conclusions}\label{sec:conclusions}
We have generalised bounds for the volume of tubular neighbourhoods from the algebraic setting~\cite{lotz2015volume} to the case of smooth Pfaffian hypersurfaces (Theorem~\ref{thm:prob_bound}), and applied these to obtain condition number tail bounds for neural network classifiers with Pfaffian activations (Section~\ref{sec:robustness}). While these results are of theoretical interest, they point to two natural directions for further development.

\subsection{Improved bounds for neural networks}\label{subsec:improved-nn}
The constant $C_{\alpha,\beta,s,n}$ in Theorem~\ref{thm:prob_bound} contains the factor $2^{s(s-1)/2}$, which for a sigmoid network with $h$ hidden units becomes $2^{h(h-1)/2}$ (Example~\ref{ex:sigmoid-condition}). The factor arises in the Khovanskii induction (Theorem~\ref{thm:khovanskii}), which peels chain elements one at a time, introducing one doubling per element and $\sum_{j=1}^{h-1}j = h(h-1)/2$ doublings in total.

The Pfaffian chain of a sigmoid network, however, has \emph{layered structure} that a generic chain does not: the chain elements decompose as $\boldq = (\boldq^{(1)}, \ldots, \boldq^{(\ell)})$, where $\boldq^{(i)} = (\sigma(z_1^{(i)}), \ldots, \sigma(z_{w_i}^{(i)}))$ are the activations of layer~$i$. This layered chain has two key properties. First, \emph{intra-layer independence:} the chain polynomial of $q_k^{(i)}$ depends on $q_k^{(i)}$ itself and on elements of earlier layers, but not on any other element $q_{k'}^{(i)}$ in the same layer. Second, \emph{layer-dependent degree:} the chain polynomial of a layer-$i$ element has degree $2i$, not the global maximum $2\ell$.

These properties yield at most polynomial improvements within the standard element-by-element induction, because the exponential factor is determined by the \emph{number} of Rolle steps, not by the per-step degree. An exponential improvement would require processing entire layers simultaneously. The key observation is that the layer map $\boldq^{(\ell-1)} \mapsto \boldq^{(\ell)}$ factors as $\sigma \circ (\text{affine})$, with Jacobian $\diag(\sigma'(z_1), \ldots, \sigma'(z_w)) \cdot A^{(\ell)}$, where $\sigma'(z_k) = q_k^{(\ell)}(1 - q_k^{(\ell)}) > 0$ for all finite $z_k$. The strict positivity of $\sigma'$ and the intra-layer independence of the chain polynomials are structural features that a \emph{block} analogue of the Khovanskii--Rolle lemma could exploit.

\begin{conjecture}[Block Khovanskii bound]\label{conj:block-khovanskii}
  Let $f\colon \R^n \to \R$ be the output of a fully connected neural network with $\ell$ hidden layers of widths $w_1, \ldots, w_\ell$ and a Pfaffian activation function $\sigma$ satisfying $\sigma' > 0$. Then the number of regular solutions of a Pfaffian system involving $f$ and its derivatives, of the form~\eqref{eq:pfaffian-system}, is bounded by
  \begin{equation*}
    2^{\frac{\ell(\ell-1)}{2} + \sum_{i=1}^{\ell}\frac{w_i(w_i-1)}{2}} \cdot p(\alpha, \beta, n, \ell, w_1, \ldots, w_\ell),
  \end{equation*}
  where $p$ is a polynomial in the indicated quantities.
\end{conjecture}

For constant width $w$, the conjectured exponent is $\frac{\ell(\ell-1)}{2} + \frac{\ell w(w-1)}{2}$, which is linear in both $\ell$ and $w$, compared with the standard $\frac{\ell w(\ell w-1)}{2} \approx \frac{\ell^2 w^2}{2}$. The improvement is substantial: for a network with $\ell = 10$ layers of width $w = 5$, the exponent drops from $1225$ to $145$, a factor of $2^{1080}$.

It is worth emphasising how Conjecture~\ref{conj:block-khovanskii} relates to Theorem~\ref{thm:multi-layer-degree}. The latter already achieves a \emph{fully polynomial} bound $\md(V_{ij})\leq C(n)\cdot (\prod_i w_i)^n$ for sigmoid classifier decision boundaries, eliminating the Khovanski\u i exponential entirely. However, that result is specific to the Gauss-map system arising from a classifier, and relies on the activation-coordinate change of Section~\ref{subsec:multi-layer}, which is available only when the ambient equation system admits an $n$-coordinate parametrisation by activation values. Conjecture~\ref{conj:block-khovanskii} is intended as an intermediate statement that applies to \emph{arbitrary} Pfaffian systems involving a layered sigmoid chain (e.g., general Pfaffian complete intersections in the sense of Proposition~\ref{prop:pfaffian-degree-bound}, not only Gauss-map systems of classifiers). For the specific case of classifier decision boundaries, Theorem~\ref{thm:multi-layer-degree} is strictly sharper; the conjecture is motivated by settings where the activation-coordinate reduction does not apply and one still wishes to avoid the full Khovanski\u i doubling.

The main obstacles are twofold. First, the standard Rolle lemma is inherently one-dimensional: it uses the intermediate value theorem on arcs of a $1$-manifold, and multi-dimensional analogues that give effective, multiplicity-free bounds remain open. Second, even with intra-layer independence, degrees compound across within-layer Rolle steps in the standard sequential argument. A proof of Conjecture~\ref{conj:block-khovanskii} would require either a multi-dimensional Rolle lemma or a refined degree-tracking scheme that separates within-layer and between-layer contributions.

\subsection{Non-smooth Pfaffian sets}\label{subsec:nonsmooth}
Theorem~\ref{thm:prob_bound} requires the Pfaffian hypersurface $V$ to be smooth, in the sense that $\nabla f$ is non-vanishing on $V$. This is a genuine restriction: the decision boundaries of neural network classifiers can develop singularities for certain weight configurations, and the volume bounds of Section~\ref{sec:tubular} do not apply at such points.

In the algebraic setting, Basu and Lerario~\cite{basu2021hausdorff} extended tube volume bounds to singular varieties by approximating a singular algebraic set $Z$ with a family $\{Z_t\}_{t > 0}$ of smooth complete intersections converging to $Z$ in the Hausdorff metric, and showing that the tube bounds pass to the limit. Their construction uses deformations of the form $D(Q, G, \zeta) = (1-\zeta)Q - \zeta G$, where $Q$ defines the original variety and $G$ is a polynomial built from Morse functions. A crucial step is to show that the set of ``bad'' parameter values---those for which the deformed variety fails to be a smooth complete intersection---forms a proper \emph{Zariski-closed} subset of $\C$. This relies on Bertini-type theorems over $\C$ and the fact that polynomial equations define constructible sets in the Zariski topology.

The fundamental obstacle to extending this strategy to Pfaffian functions is that the Zariski-constructibility arguments are intrinsically complex-algebraic. Pfaffian functions are real-analytic and have no natural complex extension that preserves their Pfaffian structure; in particular, there is no Pfaffian analogue of the fact that the singular locus of a complex algebraic variety is itself algebraic. As Basu and Lerario note~\cite[Remark~1.2]{basu2021hausdorff}, ``it is not clear if our technique can be extended to the definable setting, the main obstacle being the extension of the definition of polar varieties and their properties coming from complex algebraic geometry.''

Several avenues for progress can be envisaged.
\begin{enumerate}
  \item \emph{O-minimal Hausdorff limits.} The convergence argument in~\cite{basu2021hausdorff} can potentially be recast in o-minimal terms, since Pfaffian sets generate an o-minimal structure~\cite{wilkie1999theorem} in which Hausdorff limits of definable families are definable. The challenge is to obtain explicit format bounds for the limiting set.
  \item \emph{Pfaffian polar varieties.} Developing a theory of polar varieties for Pfaffian functions---with effective bounds on the Pfaffian format of the polar loci---would provide the geometric control needed for deformation arguments. The main risk is that generic transversality results, which are immediate in the algebraic setting, may require new techniques in the Pfaffian context.
  \item \emph{Multidimensional variations.} An alternative approach, based on Vitushkin's theory and developed in the definable setting by Comte and Yomdin, bounds tube volumes via multidimensional variations rather than degree. This approach applies in principle to Pfaffian sets but yields bounds with exponential dependence on the ambient dimension $n$~\cite[Remark~1.2]{basu2021hausdorff}, which may be acceptable in low-dimensional settings.
\end{enumerate}

